\chapter{Scripts MATLAB e Simulink}
\label{ap:matlab}

Todos os arquivos ficam em \code{guia\_controle/matlab/}. Abra o MATLAB nessa pasta.
Requisitos: MATLAB R2019b ou mais novo e \textbf{Control System Toolbox}.
\textbf{Opcionais}: Optimization (\code{quadprog}), Model Predictive Control
(\code{mpc}), System Identification (\code{tfest}), Statistics (\code{kruskalwallis})
e Simulink. Os scripts testam se cada toolbox opcional existe antes de usá-la.

\section{Organização}
\begin{table}[H]
\centering
\small
\begin{tabular}{@{}lp{0.62\textwidth}@{}}
\toprule
Arquivo & Função \\
\midrule
\code{parametros.m} & Todos os parâmetros (x500, PX4, atrasos, KF, ganhos do \code{yaml}) \\
\code{modelo\_px4.m} & $T_v(s)$, plantas $P_a$, $P_v$ contínuas e discretas (\cref{cap:discreto}) \\
\code{analise\_margens.m} & \Cref{tab:margens} e \cref{fig:margem-atraso,fig:px4} \\
\code{projeto\_pid.m} & Alocação, LGR discreto, Ziegler--Nichols, PI ótimo LQ (\cref{cap:pid}) \\
\code{projeto\_lqr.m} & Fórmulas analíticas, Bryson, LGR simétrico, margens, preditor (\cref{cap:lqr}) \\
\code{projeto\_lqi.m} & Escolha de $q_\xi$, rejeição de vento (\cref{cap:lqi}) \\
\code{projeto\_smc.m} & $\lambda^\star$, PD equivalente, camada-limite, chattering (\cref{cap:smc}) \\
\code{projeto\_mpc.m} & $K_{MPC}$, horizonte, peso terminal = DARE, \code{quadprog}, MPC Toolbox (\cref{cap:mpc}) \\
\code{projeto\_mpc\_dinamico.m} & Modelo de corpo rígido x modelo com atitude da PX4 (\cref{cap:mpcdin}) \\
\code{projeto\_hermite.m} & $T^\star$, ganho equivalente não linear (\cref{cap:hermite}) \\
\code{projeto\_guiamento.m} & Malha radial, varredura de $N$ e $\alpha$ (\cref{cap:guiamento}) \\
\code{identificar\_px4.m} & Identificação da malha de velocidade a partir de log \\
\code{cenario.m} & Trajetórias do alvo da comparação \\
\code{sim\_perseguicao.m} & Simulação 3D de malha fechada: PX4 + câmera + atraso + ruído + KF \\
\code{otimizar\_ganhos.m} & Ajuste ótimo de qualquer controlador pelo custo comum \eqref{eq:jcomum} \\
\code{comparar\_controladores.m} & Monte Carlo, tabela, boxplots, Pareto, testes estatísticos \\
\code{controladores/ctrl\_*.m} & As 7 leis, com as mesmas equações do C++ \\
\code{simulink/criar\_modelo\_simulink.m} & Gera \code{perseguicao\_1eixo.slx} por código \\
\code{simulink/rodar\_simulink.m} & Simula PID, LQR, LQI e SMC no modelo Simulink \\
\bottomrule
\end{tabular}
\end{table}

\section{Interface comum dos controladores}
Cada lei é uma função
\begin{lstlisting}
[cmd, mem] = ctrl_xxx(st, mem, g, P)
% st : p_rel, v_rel, v_chaser (= v_ref), v_target, dt   (3x1, ENU)
% mem: memoria interna (integradores, warm start); [] na primeira chamada
% g  : ganhos (ex.: P.lqr)        P: parametros()
% cmd: v (3x1), a (3x1)
\end{lstlisting}
igual a \code{PursuitController::compute} no C++. Para testar uma lei nova, basta
criar \code{ctrl\_nova.m} com essa assinatura e passá-la a \code{sim\_perseguicao}.

\section{Fluxo sugerido de estudo}
\begin{enumerate}
  \item \code{analise\_margens} --- entender a planta (PX4 + atraso).
  \item \code{projeto\_lqr} --- a referência ótima; depois \code{projeto\_pid},
        \code{projeto\_lqi}, \code{projeto\_smc}, \code{projeto\_mpc} etc.
  \item \code{simulink/rodar\_simulink} --- a mesma comparação em diagrama de blocos,
        com um eixo. É útil para experimentar atraso, amostragem e PX4 mais lenta.
  \item \code{G.lqr = otimizar\_ganhos('lqr'); ...; save ganhos\_otimos.mat G}.
  \item \code{comparar\_controladores} --- resultados para o texto.
  \item Copie os ganhos para o \code{aero\_pursuit.yaml} e repita no Gazebo.
\end{enumerate}

\section{Equivalência com o C++}
\code{sim\_perseguicao.m} reproduz o \emph{harness} C++ usado no \cref{cap:caso}:
integração a \SI{2}{ms}, PX4 com $K_{P,v}$, $K_{I,v}$, $K_{D,v}$ e atitude de 1ª ordem,
câmera a \SI{30}{Hz} com atraso e ruído, KF com \emph{gating} e reinício após 5 rejeições,
controlador a \SI{20}{Hz} e \code{applyLimits}. Para reproduzir o estimador de
\emph{antes} da correção (causa 1 do \cref{cap:caso}), use \code{P.kf\_compensar = false}
e \code{P.kf\_sigma\_z = 0.05}. A causa 2 (referência de velocidade do MPC) já está
corrigida em \code{ctrl\_mpc.m}.
